\chapter{Conclusion}\label{ch:conclusion}
\graphicspath{{../Figures/80-conclusion/}{figs/80-conclusion/}}

This dissertation moved the biomimetic humanoid robot project from an unvalidated design concept to a quantitatively characterized actuator set and experimentally tested knee-torque models, and it defined the path from those experimental foundations to neuromechanical control.

At the actuator level, the analysis in Chapter~\ref{ch:muscles} produced novel equations for calculating isometric force in Festo \qtylist{10;20}{\mm}-diameter BPAs. The work characterized the relationship between maximum BPA force at \qty{620}{\kPa} and resting length, $F_{620}(l_0)$, a previously unreported characteristic of these actuators that the manufacturer's tool predicts with the opposite trend, and it produced a nondimensional force surface, $F^*(\epsilon^*,P^*)$, with only three diameter-dependent coefficients. Taken together, these fits predict isometric BPA force across resting lengths from a small number of measurements, correct the large over-predictions of the manufacturer's tool at short resting lengths, and outperform published single-length models. In a design workflow, these equations make it possible to select actuator diameter and resting length and to evaluate whether a proposed muscle routing can meet force and range-of-motion requirements, before committing to a prototype.

At the joint level (Chapter~\ref{ch:torque}), the BPA model was combined with transmission corrections for effective length, series compliance, and wrapping. Identification on the pinned knee improved torque prediction, and the shared correction terms also improved predictions on the biomimetic four-bar knee. The corrected model supports route redesign: the optimized two-BPA flexor and cam-extensor configurations are predicted to meet their human-scaled torque targets across the design range. The isometric tests demonstrate that the redesigned configurations meet or exceed their selected human-muscle reference torques at all fifteen flexor and nine extensor measurement points. This establishes capacity at the sampled angles and pressures, while agreement with pressure-specific predictions remains uneven.

The SADb application in Chapter~\ref{ch:database} makes physiological evidence searchable by pathway, preparation, and modeling relevance. The tools in Chapter~\ref{ch:sns} turn selected pathways into testable neural--mechanical models: the Li--Deng/W2L and Gait2392 implementations, SNS-Simscape blocks, MuJoCo--Simulink integration, and staged tuning and intervention procedures. Demonstrations establish matched-circuit agreement, corrective feedback, command and measurement exchange, and pathway-dependent stepping under specified support. These tests provide a repeatable framework for comparing neural and mechanical hypotheses. General controller transfer and independent balance remain open.

These actuator and joint models, the sensory-afferent application, and the neuromechanical tools make design assumptions testable: force against measurement, routing against torque requirements, and sensory pathways against controlled changes in movement. The future studies in Chapter~\ref{ch:overall} extend these tests toward dynamic actuation, independent balance, and hardware. This work provides an experimental foundation for biomimetic robot design and for studying how neural control interacts with mechanical structure.
